\begin{circuitikz}[
    american,
    voltage dir=RP,
    >=latex,
    %>={Latex[length=5mm, width=2mm]},
    alt={}
    ]

    % Define the number of rungs in the diagram
    \def\numRungs{2}

    % Calculate the length of the vertical power rails dynamically based on number of rungs
    \pgfmathsetmacro{\railLength}{-2 * \numRungs}

    %======================================================================================================
    % Component Grid
    % Spacing = 3x, 2y
    %======================================================================================================
    \foreach \i in {0,...,5} {
        \foreach \j in {0,...,20} { % Change to 0,...,19 if you need 20 individual Y points down to -40
        
        % Calculate spatial positions
        \pgfmathsetmacro{\px}{\i * 3}
        \pgfmathsetmacro{\py}{-1 - (\j * 2)}
        
        % Name coordinate using loop indices: (C-column-row) e.g., (C-0-0), (C-5-10)
        \coordinate (C-\i-\j) at (\px, \py);
        
        % OPTIONAL: Uncomment the 2 lines below to display grid labels while building
        % \fill[red] (C-\i-\j) circle (1.5pt);
        %\node[anchor=south west, scale=0.5, gray] at (C-\i-\j) {(\i,\j)};
        }
    }

    %======================================================================================================
    % Foreground Layer
    % Only contains the vertical power rails
    % Length of the rails is automatic, calculated based on the number of rungs defined above
    %======================================================================================================
    \begin{pgfonlayer}{ft}

        %======================================================================================================
        % Scope sets options for both lines - BakerBlack, very thick
        %======================================================================================================
        \begin{scope}[
            draw=BakerBlack,
            text=BakerBlack,
            color=BakerBlack,
            font=\small,
            ultra thick
        ]
        
            \draw (0,0) -- (0,\railLength);

            \draw (15,0) -- (15,\railLength);

        \end{scope}
    \end{pgfonlayer}

    %======================================================================================================
    % Main Layer
    % All actual logic goes here
    %======================================================================================================
    \begin{pgfonlayer}{main}
        
        %======================================================================================================
        % Scope sets options for both lines - BakerBlack for draw and text
        %======================================================================================================
        \begin{scope}[
            draw=BakerBlack,
            text=BakerBlack,
            color=BakerBlack,
            font=\small
        ]

            %=================================================================
            % Rung 1
            %=================================================================
            \draw (C-0-0) -- (C-3-0) pic (TON) {TON};
            \draw (TON-OUT) -- (C-5-0);

            \node[anchor=south] at (TON-TOP) {TON};
            \node[anchor=north] at (TON-TOP) {\texttt{RunTimer}};
            \node[anchor=west] at (TON-PRE) {20000};
            \node[anchor=west] at (TON-ACC) {6789};

            \draw ($(C-2-0)+(2,0)$) coordinate(h1) -- ++(0,-3) coordinate(h2) -- ++(0.5,0) pic (SUB) {SrcDest};
            \draw (SUB-OUT) -- ++(0.5,0) coordinate(h3) -- ++(0,3) coordinate(h4);

            \node[anchor=south] at (SUB-TOP) {SUB};
            \node[anchor=west] at (SUB-SRCA) {\texttt{RunTimer.PRE}};
            \node[anchor=west] at (SUB-AVAL) {20000};
            \node[anchor=west] at (SUB-SRCB) {\texttt{RunTimer.ACC}};
            \node[anchor=west] at (SUB-BVAL) {6789};
            \node[anchor=west] at (SUB-DEST) {\texttt{TimeLeft\_ms}};
            \node[anchor=west] at (SUB-DVAL) {13211};

            \draw (h2) -- ++(0,-4.75) coordinate(h5) -- ++(0.5,0) pic (DIV) {SrcDest};
            \draw (DIV-OUT) -- ++(0.5,0) coordinate(h6) -- ++(0,4.75) coordinate(h7);

            \node[anchor=south] at (DIV-TOP) {DIV};
            \node[anchor=west] at (DIV-SRCA) {\texttt{TimeLeft\_ms}};
            \node[anchor=west] at (DIV-AVAL) {13211};
            \node[anchor=west] at (DIV-SRCB) {1000};
            \node[anchor=west] at (DIV-BVAL) {1000};
            \node[anchor=west] at (DIV-DEST) {\texttt{TimeLeft\_S}};
            \node[anchor=west] at (DIV-DVAL) {13.211};

        \end{scope}
    \end{pgfonlayer}

    %=================================================================
    % Background layer
    %=================================================================
    % \begin{pgfonlayer}{bg}
        
    %     %==================================================================
    %     % Highlight Scope
    %     % This layer exists for highlighting true logic, the scope sets the
    %     % highlight options without having to repeat them every \draw
    %     % draw=UpNorthGreen, line width=10pt, opacity=0.4
    %     %==================================================================
    %     \begin{scope}[
    %         draw=UpNorthGreen,
    %         line width=10pt,
    %         opacity=0.4
    %     ]

    %         %==================================================================
    %         % Left vertical rung - leave this alone, it should ALWAYS be green
    %         % It automatically draws itself to the length defined at the top
    %         %==================================================================
    %         \draw (0,0) -- (0,\railLength);
        
    %         %=================================================================
    %         % Rung 1
    %         %=================================================================
    %         \draw (C-0-0) -- (TON-IN);
    %         % \draw (MOV1-OUT) -- (C-5-0);

    %         %\draw (h1) -- (h2) -- ({SUB-IN});
    %         % \draw (CPT-OUT) -- (h3) -- (h4);

    %     \end{scope}
    % \end{pgfonlayer}

\end{circuitikz}